\section{Heights, two spins, and exact cuts}\label{sec:spin}

Our first objective is to replace the height constraint by a finite spin
law whose conditional distributions can be read exactly.  The loop
description will connect this law to the scale estimates in
\cref{sec:cuts}; the one-color cuts will remain available when rare
boundary pins are imposed later.
We use the coherent-pair and double-path framework of
Glazman--Manolescu \cite[Proposition~1.4 and Sections~2.1--2.3]{GM2021},
and establish the versions with equality constraints and stopping cuts
needed below.

We use lattice units in the discrete arguments.  A finite triangular domain
is a finite, simply connected triangular cell complex, with its vertex and
edge sets understood.  A polygonal boundary, or a separating lattice curve,
includes the adjacent fixed vertices when a Gibbs interior is taken.  In
particular, a free vertex has its full incident star up to fixed boundary
vertices.  This convention avoids an additional interaction across the
boundary of a conditional problem.

Write $H=(h-1)/2$.  Thus $H$ is integer-valued and
$\abs{H(x)-H(y)}\leq1$ on each edge.  Associate two spins to a height by
\begin{equation}\label{eq:gray-code}
 B=+1\quad\hbox{on residues }0,1,\qquad
 W=+1\quad\hbox{on residues }-1,0\pmod4.
\end{equation}
A pair $(B,W)$ is \emph{coherent} if the two spins never both change across
one edge.  A free spin law means uniform counting of coherent pairs,
without boundary conditions.

\begin{lemma}[The lift and the loop representation]\label{lem:gray}
On a simply connected triangular domain, a coherent pair determines
integer height increments
\begin{equation}\label{eq:height-increment}
 dH(x,y)=\frac{W_xB_y-B_xW_y}{2}.
\end{equation}
They have zero curl.  Fixing the height at one vertex gives a bijection
between coherent pairs whose Gray state there matches that prescribed
integer and admissible heights with that value.  Uniform counting has the
height Gibbs property.

With a constant boundary height, the nonzero height steps form disjoint
closed honeycomb loops.  Conditional on their geometry, the inward jumps
are independent fair signs; the unoriented loop weight is
$2^{\#\mathrm{loops}}$.  The same conclusion holds in a bounded face of a
connected skeleton of fixed lattice edges carrying one constant height.
\end{lemma}

\begin{proof}
The four Gray states, in increasing residue order, are
$(+,+),(+,-),(-,-),(-,+)$.  Inspection of consecutive states gives
\eqref{eq:height-increment}.  On a triangle, if $B$ is not constant, it
changes on two incident edges; coherence forces $W$ to be constant on all
three vertices.  Interchanging the colors gives the alternative case.
The sum of \eqref{eq:height-increment} around that triangle is therefore
zero.  Every closed path is generated by triangle boundaries, proving
path independence and the asserted bijection.  Given a separating height
trace, all compatible interiors have the same weight, which proves the
Gibbs assertion.

A triangle has either no nonzero step or two nonzero steps.  In the latter
case its three heights take two consecutive values, so the two associated
dual edges join into one contour separating those same levels.  Contours
cannot branch.  A constant boundary has no contour endpoints, and each
loop may independently be assigned an inward jump $+1$ or $-1$.
Summing these jumps from the boundary constructs the inverse height
configuration.  Thus every unoriented loop configuration has exactly
$2^{\#\mathrm{loops}}$ compatible orientations.  A loop in a bounded
complementary face cannot enclose only part of a connected fixed skeleton:
a path in the skeleton from that part to the rest would cross the loop.
The same construction consequently applies in such a face.
\end{proof}

The spin comparison must tolerate more than constant boundary conditions.
In a finite free domain let $E_0$ be any set of edges on which $W$ is required
to be constant, and prescribe arbitrary $B$ values on an arbitrary vertex
set.  Write $\mu^{E_0,\eta}$ for this constrained pair law.  An empty set of
$B$ prescriptions is allowed.  These constraints are always feasible: take
$W$ constant and then assign $B$ arbitrarily.

\begin{lemma}[Spin FKG with equality constraints]\label{lem:spin-fkg}
The $B$ marginal of $\mu^{E_0,\eta}$ satisfies the FKG lattice condition.
The assertion remains true after any further $B$ pinnings.  In particular,
fixing additional $B$ spins to $+1$ increases increasing events, and fixing
them to $-1$ decreases such events.

Conditional on $B$, the spin $W$ consists of independent fair signs on the
components of the graph generated by $E_0$ and the $B$-changing edges.
\end{lemma}

\begin{proof}
Encode a $W$ wall by a binary variable on each edge.  Such variables are
the differences of vertex signs if and only if their sum around every
triangle is zero over $\mathbb F_2$.  Simple connectivity makes these
triangle equations sufficient; each solution has two vertex lifts.
Coherence and $E_0$ require the wall variable to vanish except on
equal-$B$ edges outside $E_0$.

Let $A$ be the fixed triangle--edge incidence matrix over $\mathbb F_2$.
For a set $F$ of allowed columns put
$n(F)=\abs F-\operatorname{rank}(A_F)$.  This nullity is increasing and
supermodular.  Indeed, adding one column increases rank by at most one,
and submodularity of matrix rank gives
\[
 n(F\cap G)+n(F\cup G)\geq n(F)+n(G).
\]
Let $F_+(B)$ and $F_-(B)$ be the allowed edges whose two $B$ endpoints
are respectively $+1$ and $-1$.  A triangle cannot contain both a $++$
edge and a $--$ edge.  Its parity equations thus split into two disjoint
systems, and the number of compatible $W$ configurations is
\begin{equation}\label{eq:spin-nullity-weight}
 w(B)=2^{1+n(F_+(B))+n(F_-(B))}.
\end{equation}
For the plus edge set,
\[
 F_+(B\wedge B')=F_+(B)\cap F_+(B'),\qquad
 F_+(B\vee B')\supseteq F_+(B)\cup F_+(B').
\]
Monotonicity and supermodularity of $n$ give its contribution to
$w(B\wedge B')w(B\vee B')\geq w(B)w(B')$.  For the minus edge set the
same argument has meet and join interchanged.  Multiplication proves the
lattice condition.  An indicator of fixed coordinate values preserves
this condition, including when one of the compared weights is zero.
The standard FKG conditional comparison gives the stated monotonicity.

Alternatively describing $W$ by its vertices, coherence forces equality
exactly across $B$-changing edges, in addition to the equalities in $E_0$.
There are no other restrictions.  Uniform counting gives one independent
fair sign per resulting component.
\end{proof}

We also need association and boundary comparison for heights.
Coordinatewise minimum and maximum of two integer Lipschitz functions
are integer Lipschitz.  If their boundary values are ordered, sorting by
minimum and maximum preserves the respective boundary conditions.
This is the finite-lattice setting of
Fortuin--Kasteleyn--Ginibre \cite{FKG1971}.  A fixed boundary or anchor
bounds every height by graph distance, so the height coordinates lie in
finite chains.

The following finite-state argument proves height association and
monotonicity in the boundary values, as well as the spin comparison
consequences of \cref{lem:spin-fkg}.  Update one unfixed coordinate at a time by its
conditional quantile, using independent uniforms and a deterministic
cyclic schedule.  For heights its conditional support is the integer
interval
\[
 [\max_{w\sim v}H(w)-1,\ \min_{w\sim v}H(w)+1],
\]
intersected with the fixed coordinate bounds; both endpoints are
increasing in the neighboring heights.  For the first spin the lattice
condition makes its conditional plus probability increasing in all
other unfixed spins.  Thus each update is increasing in the current
configuration and in its uniform random variable.  Starting from the
least feasible configuration, each finite sequence of updates is an
increasing function of independent variables and hence is positively
associated, by the one-dimensional Chebyshev inequality and induction on
the independent coordinates.  The same uniforms couple ordered boundary prescriptions
in their indicated order.

These finite chains converge to their counting laws.  The spin marginal
is positive on every assignment of the unpinned first-color spins,
because a constant second color is always coherent.  For heights, any
configuration can be decreased to the coordinatewise minimum with a
second one: lower a discrepant vertex of largest current height by one.
A higher neighbor already at its target would contradict the target's
Lipschitz condition, so this move is feasible.  Repeating and reversing
the corresponding moves for the second configuration proves
irreducibility.  Each single-coordinate conditional law has a positive
holding probability.  A cyclic sweep therefore gives an irreducible
aperiodic chain preserving the desired law.  Passing to its stationary
limit proves association and the ordered couplings, including the
specified first-color pinnings and second-color equalities.

\begin{definition}[Double paths]
A dual edge is \emph{double-$B$ of sign $t$} if the two primal endpoints
of the crossed edge both have $B=t$.  A double path or circuit uses only
such edges, with one fixed sign throughout.  The analogous terminology
applies with the colors interchanged.
\end{definition}

\begin{lemma}[Exact cut law]\label{lem:double-cut}
Let $C$ be a simple double-$B$ circuit of sign $t$.  Suppose that no edge
of $E_0$ crosses $C$.  Conditional on the exterior pair configuration and
on the $B=t$ values on the inner boundary layer of $C$, the interior law
is the uniform coherent-pair law with that $B$ boundary layer, free $W$
boundary values, and just the constraints situated inside $C$.

This assertion also holds for a circuit selected by a stopping exploration
which reveals only exterior $B$ spins and the $B=t$ cut sites.  No $W$
spins need be revealed during that exploration.  Along the inner boundary
face chain, a height lift takes values in one pair of consecutive integers.
\end{lemma}

\begin{proof}
Every interaction crossing $C$ is on an edge with $B$ equal at both
endpoints.  Coherence places no restriction on $W$ across such an edge,
and the hypothesis excludes an $E_0$ equality there.  Once the cut $B$
values are fixed, the counting weight therefore factors into an exterior
factor and the stated interior factor.  Fixed constraints inside and
outside occur in their respective factors.

For the stopping version, explore from the designated exterior side
through primal edges not blocked by a good dual edge, querying the two
endpoint $B$ values when needed.  If a surrounding double circuit exists
in the search collar, this exploration cannot pass it.  The separating
boundary contains a simple double circuit.  The queried sites on its
inner side are only its good cut sites; bounded pockets may be filled on
the already explored side.  Thus, conditional on the exploration record,
the same factorization applies.  This is a finite exploration statement,
so it also applies after arbitrary exterior conditioning of the allowed
type.

At a dual vertex of a double path, its two used edges force all three
adjacent $B$ spins to equal $t$.  The inner face sites of a double circuit
therefore form a connected constant-$B$ chain.  The two residues with this
$B$ value are consecutive.  An integer Lipschitz lift along a connected
constant-$B$ chain stays in one translated copy of that pair, proving the
last assertion.
\end{proof}

We shall impose $W=+1$ along a bounded number of connected paths.  It is
useful to separate their equality constraints from their labels.  First
put their edges in $E_0$.  Conditional on $B$, each path then has one
component label.  If there are $m$ paths, the conditional probability of
all labels being $+1$ lies between $2^{-m}$ and $1$.  If there is just one
connected path, this probability is exactly $1/2$, independently of $B$;
requiring its plus label does not change the $B$ marginal.  We will always
account for these labels after an argument involving the $B$ marginal.
